% ============================================================
%  CHAPTER 24: ELECTRICAL ENERGY AND POWER
%  The Physics Codex: A Complete University Foundation
%  Korede Samuel Omotosho (Falcon)
%  Aurikrex Learning Systems, 2026
%
%  File: chapters/part4/ch24_electrical_power.tex
% ============================================================

\chapter{Electrical Energy and Power}
\label{ch:electrical_power}
\minitoc
\newpage

\chapterquote{Energy is the only life, and is from the
Body; and Reason is the bound or outward circumference
of Energy.}{William Blake}

\begin{objectivebox}
By the end of this chapter, you should be able to:
\begin{enumerate}[noitemsep]
  \item Define electrical power and calculate it
  using multiple formulae
  \item Calculate electrical energy in joules and
  kilowatt-hours
  \item Explain the heating effect of electric
  current and apply it to practical situations
  \item Describe the action of fuses and circuit
  breakers
  \item Explain why electrical power is transmitted
  at high voltage and low current
  \item Calculate power losses in transmission lines
  \item Describe the principles of AC generation
  and the transformer
  \item Apply the transformer equations for voltage,
  current and power
  \item Perform cost calculations for electrical
  energy consumption
\end{enumerate}
\end{objectivebox}

% ============================================================
\section{Intuition: Energy from Electrons}
% ============================================================

When you switch on a $60\ \text{W}$ light bulb, you
are driving electrons through the tungsten filament
at roughly $0.27\ \text{A}$. Each electron loses
$222\ \text{V}$ worth of potential energy as it moves
through the filament. That energy does not disappear ---
it is converted to heat and light, heating the filament
to about $2500°\text{C}$, hot enough to glow white.

The rate at which electrical energy is converted to
other forms (heat, light, mechanical work) is
\textbf{electrical power}. It is simply the product
of voltage and current. This chapter connects the
abstract quantities of current and potential difference
to the everyday reality of electricity bills, burnt-out
fuses, power transmission lines, and the transformer
sitting on every electrical pole across Nigeria.

% ============================================================
\section{Electrical Power}
% ============================================================

\begin{definitionbox}[Electrical Power]
The \textbf{electrical power} dissipated in a
component is the rate at which electrical energy
is converted to other forms:
\[
  P = \frac{W}{t} = \frac{QV}{t} = IV
\]
where $W = QV$ is the work done (energy transferred)
when charge $Q$ moves through potential difference $V$.

Using Ohm's Law ($V = IR$), three equivalent forms:
\[
  \boxed{P = IV = I^2R = \frac{V^2}{R}}
\]
SI unit: watt (W) where $1\ \text{W}
= 1\ \text{J\,s}^{-1}$.

\textbf{Which formula to use:}
\begin{itemize}[noitemsep]
  \item Know $I$ and $V$: use $P = IV$
  \item Know $I$ and $R$: use $P = I^2R$
  \item Know $V$ and $R$: use $P = V^2/R$
\end{itemize}
\end{definitionbox}

\begin{examplebox}[24.1]
\stepcounter{workedexample}
An electric iron is rated $1200\ \text{W}$ at
$240\ \text{V}$.\\
(a) Calculate the current drawn.\\
(b) Calculate the resistance of the heating element.\\
(c) Calculate the charge passing through the iron
in $30\ \text{min}$.\\
(d) Calculate the energy consumed in $30\ \text{min}$.
\end{examplebox}

\begin{solutionbox}
\textbf{(a)} Current:
\[
  I = \frac{P}{V} = \frac{1200}{240} = 5.0\ \text{A}
\]

\textbf{(b)} Resistance:
\[
  R = \frac{V}{I} = \frac{240}{5.0} = 48\ \Omega
\]
Or: $R = V^2/P = 240^2/1200 = 57600/1200 = 48\ \Omega$ ✓

\textbf{(c)} Charge in $30\ \text{min}
= 1800\ \text{s}$:
\[
  Q = It = 5.0 \times 1800 = 9000\ \text{C}
\]

\textbf{(d)} Energy consumed:
\[
  W = Pt = 1200 \times 1800 = 2\,160\,000\ \text{J}
  = 2.16\ \text{MJ}
\]
\end{solutionbox}

% ============================================================
\section{Electrical Energy and the Kilowatt-Hour}
% ============================================================

\begin{definitionbox}[Electrical Energy and the
Kilowatt-Hour]
Electrical energy consumed:
\[
  W = Pt \quad [\text{joules, when } P \text{ in
  watts and } t \text{ in seconds}]
\]

For billing purposes, the unit used is the
\textbf{kilowatt-hour} (kWh):
\[
  1\ \text{kWh} = 1000\ \text{W} \times 3600\ \text{s}
  = 3.6\times10^6\ \text{J} = 3.6\ \text{MJ}
\]
\textbf{Energy in kWh:}
\[
  W(\text{kWh}) = P(\text{kW}) \times t(\text{hours})
\]
\textbf{Cost of electrical energy:}
\[
  \text{Cost} = W(\text{kWh}) \times \text{tariff rate}
  (\text{per kWh})
\]
\end{definitionbox}

\begin{examplebox}[24.2]
\stepcounter{workedexample}
A household uses the following appliances daily in
Nigeria:
\begin{itemize}[noitemsep]
  \item One $1.5\ \text{kW}$ air conditioner for
  $6\ \text{h}$
  \item Four $18\ \text{W}$ LED bulbs for $8\ \text{h}$
  each
  \item A $100\ \text{W}$ television for $5\ \text{h}$
  \item A $2.0\ \text{kW}$ electric iron for $1\ \text{h}$
\end{itemize}
The electricity tariff is ₦68 per kWh.\\
(a) Calculate the daily energy consumption in kWh.\\
(b) Calculate the daily cost.\\
(c) Calculate the monthly (30-day) bill.
\end{examplebox}

\begin{solutionbox}
\textbf{(a)} Daily energy:

\begin{center}
\begin{tabular}{lllll}
\toprule
Appliance & Power (kW) & Hours & kWh\\
\midrule
Air conditioner & 1.50 & 6 & 9.00\\
Four LED bulbs & $4\times0.018 = 0.072$ & 8 & 0.576\\
Television & 0.100 & 5 & 0.500\\
Electric iron & 2.00 & 1 & 2.00\\
\midrule
\textbf{Total} & & & \textbf{12.076}\\
\bottomrule
\end{tabular}
\end{center}

Daily energy $\approx 12.1\ \text{kWh}$

\textbf{(b)} Daily cost:
\[
  \text{Cost} = 12.076 \times 68
  = ₦821.17 \approx ₦821
\]

\textbf{(c)} Monthly bill (30 days):
\[
  \text{Monthly} = 821 \times 30 = ₦24\,634
\]

The air conditioner dominates (74\% of the bill).
This explains why reducing air conditioner usage
or switching to a higher-efficiency unit dramatically
reduces electricity costs.
\end{solutionbox}

% ============================================================
\section{The Heating Effect of Electric Current}
% ============================================================

When current flows through a resistor, the work done
by the electric field on the charge carriers is
converted to thermal energy through collisions with
lattice ions. This is called \textbf{Joule heating}
or the \textbf{thermal effect of current}.

\begin{lawbox}[Joule's Law of Heating]
The heat generated by a current $I$ flowing through
a resistance $R$ for time $t$ is:
\[
  Q_{heat} = I^2Rt
\]
This is always positive --- heat is always generated
regardless of the direction of current. The heating
effect is why resistors get warm, why fuses blow,
and why cables have maximum current ratings.
\end{lawbox}

\begin{historybox}[James Prescott Joule, 1840]
James Joule (1818--1889) was a Manchester brewer who
became one of the greatest experimental physicists
of the 19th century. He discovered the mechanical
equivalent of heat (that heat is a form of energy)
and the heating effect of current. He was largely
self-taught, conducting experiments in the cellar
of his family's brewery. The unit of energy (joule)
is named in his honour. His collaboration with Lord
Kelvin established the foundations of thermodynamics.
\end{historybox}

\begin{realworldbox}[Joule Heating in Nigerian Electrical Systems]
Poor quality extension cords and undersized wiring
are responsible for a significant proportion of house
fires in Nigeria. When too much current is drawn
through a wire rated for a lower current, the power
dissipated ($I^2R$) causes excessive heating. Since
resistance is fixed by the wire geometry and material,
and the heating increases as $I^2$, even a modest
overcurrent causes dramatically more heat. A wire
rated for $6\ \text{A}$ carrying $12\ \text{A}$
dissipates $4$ times the rated heat --- enough to
ignite surrounding insulation and materials.
\end{realworldbox}

% ============================================================
\section{Fuses and Circuit Breakers}
% ============================================================

\begin{definitionbox}[Fuse]
A \textbf{fuse} is a short length of thin wire
(low melting point alloy) connected in series with
a circuit. If the current exceeds the fuse rating,
the wire melts and breaks the circuit, protecting
the appliance and wiring from damage.

Fuses are rated in amperes (A): $1\ \text{A}$,
$3\ \text{A}$, $5\ \text{A}$, $13\ \text{A}$, etc.

\textbf{Correct fuse selection:} Choose the lowest
fuse rating that is above the normal operating
current of the appliance.

For appliance of power $P$ at voltage $V$:
\[
  I_{operating} = \frac{P}{V}
\]
Choose a fuse rated just above $I_{operating}$.
\end{definitionbox}

\begin{defbox}[Circuit Breaker vs Fuse]
A \textbf{circuit breaker} is an automatic switch
that trips (opens) when current exceeds a set value.
Unlike a fuse, it can be reset (switched back on)
after the fault is cleared. Modern homes use circuit
breakers (MCBs --- Miniature Circuit Breakers) in
the consumer unit (distribution board).

\textbf{Earth leakage circuit breaker} (ELCB) or
\textbf{Residual Current Device} (RCD) detects if
current is flowing to earth (through a person, for
example) and disconnects the supply within
milliseconds, preventing electrocution.
\end{defbox}

\begin{examplebox}[24.3]
\stepcounter{workedexample}
A $1000\ \text{W}$ blender and a $700\ \text{W}$
toaster are both connected to a $240\ \text{V}$
socket via a shared $13\ \text{A}$ extension cord.\\
(a) Find the operating current of each appliance.\\
(b) Find the total current drawn.\\
(c) Is the $13\ \text{A}$ fuse adequate? Explain.\\
(d) A student also plugs in a $1200\ \text{W}$
microwave oven. Will the fuse blow?
\end{examplebox}

\begin{solutionbox}
\textbf{(a)} Operating currents:
\[
  I_{blender} = \frac{1000}{240} = 4.17\ \text{A}
  \qquad
  I_{toaster} = \frac{700}{240} = 2.92\ \text{A}
\]

\textbf{(b)} Total current:
$I_{total} = 4.17 + 2.92 = 7.08\ \text{A}$

\textbf{(c)} $7.08\ \text{A} < 13\ \text{A}$,
so the fuse is adequate for this load.

\textbf{(d)} Adding microwave:
$I_{microwave} = 1200/240 = 5.00\ \text{A}$

New total: $7.08 + 5.00 = 12.08\ \text{A}$

Still just below $13\ \text{A}$, so the fuse should
not blow immediately. However, if the blender is
under heavy load and draws slightly more than rated,
or if the supply voltage fluctuates (common in
Nigeria), the total could exceed $13\ \text{A}$.
Using a $13\ \text{A}$ socket for three high-power
appliances simultaneously is poor practice.
\end{solutionbox}

% ============================================================
\section{Power Transmission and the Transformer}
% ============================================================

\subsection{Why High Voltage Transmission?}

When electrical power is transmitted over long
distances, energy is lost as heat in the cables
($P_{loss} = I^2R_{cable}$). To minimise these
losses, power is transmitted at \textbf{high voltage}
and \textbf{low current}. Since $P = IV$, the same
power can be transmitted with much lower current
if the voltage is higher --- and lower current means
much less heat dissipation.

\begin{diagbox}[High Voltage Transmission: Reducing Power Loss]
\begin{center}
\begin{tikzpicture}[scale=0.9, font=\small]

  % Power station
  \fill[gray!30] (0,1.0) rectangle (1.5,3.0);
  \draw[thick] (0,1.0) rectangle (1.5,3.0);
  \node[font=\small, align=center] at (0.75,2.0)
    {Power\\Station\\$P_{gen}$};

  % Connection to the step-up transformer
  \draw[very thick, gray] (1.5,2.0) -- (2.0,2.0);

  % Step-up transformer
  \fill[codexblue!20] (2,0.5) rectangle (3.5,3.5);
  \draw[thick] (2,0.5) rectangle (3.5,3.5);
  \node[codexblue, font=\tiny, align=center, text width=1.3cm]
    at (2.75,2.0) {Step-up\\Transformer\\$11\ \text{kV}$\\
    $\to 400\ \text{kV}$};

  % Transmission line
  \draw[very thick, gray] (3.5,2.0) -- (8,2.0);
  \draw[thick, gray, decorate,
    decoration={snake, amplitude=2pt}]
    (3.5,2.0) -- (8,2.0);
  \node[above, gray, font=\tiny, align=center]
    at (5.75,3.7) {High-voltage grid\\$400\ \text{kV}$, low current};
  \node[below, codexred, font=\small]
    at (5.75,1.85) {$P_{loss} = I^2R_{cable}$ (small)};

  % Step-down transformer (grid substation)
  \fill[codexred!20] (8,0.5) rectangle (9.5,3.5);
  \draw[thick] (8,0.5) rectangle (9.5,3.5);
  \node[codexred, font=\tiny, align=center, text width=1.3cm]
    at (8.75,2.0) {Step-down\\Transformer\\$400\ \text{kV}$\\
    $\to 11\ \text{kV}$};

  % Local distribution
  \draw[very thick, gray] (9.5,2.0) -- (12,2.0);
  \node[above, gray, font=\tiny] at (10.75,2.3)
    {$11\ \text{kV}$};

  % Local step-down
  \fill[codexgreen!20] (12,1.0) rectangle (13.5,3.0);
  \draw[thick] (12,1.0) rectangle (13.5,3.0);
  \node[codexgreen, font=\scriptsize, align=center,
    text width=1.35cm]
    at (12.75,2.0) {Local\\Transformer\\$\to 240\ \text{V}$};

  % Home
  \fill[gray!15] (14,1.2) rectangle (15,2.8);
  \draw[thick] (14,1.2) rectangle (15,2.8);
  \node[font=\scriptsize, align=center, text width=0.9cm]
    at (14.5,2.0) {Home\\$240\ \text{V}$};
  \draw[thick] (13.5,2.0) -- (14,2.0);

  % Comparison: low voltage same power
  \node[align=center, font=\small, codexblue]
    at (5.75,0.0)
    {$P = IV$: at high $V$, current $I = P/V$ is small\\
    $P_{loss} = I^2R$: small $I \Rightarrow$ tiny losses};

\end{tikzpicture}
\end{center}
\end{diagbox}

\begin{examplebox}[24.4]
\stepcounter{workedexample}
A power station generates $50\ \text{MW}$ of
electrical power. The transmission cable has a
total resistance of $10\ \Omega$.\\
(a) If transmitted at $10\ \text{kV}$, find the
current and the power lost in the cables.\\
(b) If transmitted at $400\ \text{kV}$, find the
current and the power lost.\\
(c) Compare the power losses as a percentage.
\end{examplebox}

\begin{solutionbox}
\textbf{(a)} At $10\ \text{kV}$:
\[
  I = \frac{P}{V} = \frac{50\times10^6}{10\times10^3}
  = 5000\ \text{A}
\]
\[
  P_{loss} = I^2R = (5000)^2 \times 10
  = 250\times10^6\ \text{W} = 250\ \text{MW}
\]
Power lost $= 250\ \text{MW}$ out of $50\ \text{MW}$
generated --- \textbf{5 times more power wasted than
generated}! Transmission at $10\ \text{kV}$ is
completely impractical.

\textbf{(b)} At $400\ \text{kV}$:
\[
  I = \frac{50\times10^6}{400\times10^3}
  = 125\ \text{A}
\]
\[
  P_{loss} = (125)^2 \times 10
  = 156\,250\ \text{W} = 156\ \text{kW}
\]

\textbf{(c)} Comparison:
\[
  \text{At 10 kV: } \frac{250\times10^6}{50\times10^6}
  \times 100\% = 500\%\ \text{(impossible)}
\]
\[
  \text{At 400 kV: } \frac{156\times10^3}
  {50\times10^6} \times 100\% = 0.31\%
\]
Increasing voltage by a factor of 40 reduces power
loss by a factor of $40^2 = 1600$. This is the
entire justification for the national high-voltage
grid.
\end{solutionbox}

% ============================================================
\section{The Transformer}
% ============================================================

\begin{definitionbox}[Transformer]
A \textbf{transformer} is a device that changes
the voltage of an AC supply using electromagnetic
induction. It consists of two coils (primary and
secondary) wound on a common iron core.

When AC flows in the primary coil, it creates a
changing magnetic flux in the iron core. This
changing flux induces an EMF in the secondary coil
(Faraday's Law --- Chapter 26).
\end{definitionbox}

\begin{diagbox}[Transformer Construction and Action]
\begin{center}
\begin{tikzpicture}[scale=0.95, font=\small]

  % Iron core (rectangular loop)
  \fill[gray!30] (0,0) rectangle (1.0,5.5);
  \fill[gray!30] (6.0,0) rectangle (7.0,5.5);
  \fill[gray!30] (0,5.0) rectangle (7.0,5.5);
  \fill[gray!30] (0,0) rectangle (7.0,0.5);
  \draw[very thick, gray]
    (0,0) rectangle (7.0,5.5);
  \draw[very thick, gray]
    (1.0,0.5) rectangle (6.0,5.0);
  \node[gray, font=\small] at (3.5,2.75)
    {Iron core\\(laminated)};

  % Primary coil (left leg)
  \draw[codexblue, very thick, decorate,
    decoration={coil, aspect=0.5, segment length=6pt,
    amplitude=5pt}]
    (0.5,1.0) -- (0.5,4.5);
  \node[codexblue, font=\scriptsize, align=center,
    fill=white, inner sep=1pt] at (1.7,4.0)
    {Primary\\$N_p$ turns};

  % Secondary coil (right leg)
  \draw[codexred, very thick, decorate,
    decoration={coil, aspect=0.5, segment length=6pt,
    amplitude=5pt}]
    (6.5,1.0) -- (6.5,4.5);
  \node[codexred, font=\scriptsize, align=center,
    fill=white, inner sep=1pt] at (5.3,4.0)
    {Secondary\\$N_s$ turns};

  % AC input
  \draw[thick, codexblue] (0.5,4.5) -- (-1.5,4.5);
  \draw[thick, codexblue] (0.5,1.0) -- (-1.5,1.0);
  \draw[thick, codexblue] (-1.5,4.5) -- (-1.5,2.85);
  \draw[thick, codexblue] (-1.5,1.65) -- (-1.5,1.0);
  \fill[white] (-1.5,2.25) circle (0.6);
  \draw[thick, codexblue] (-1.5,2.25) circle (0.6);
  \draw[codexblue, thick]
    (-1.9,2.25) .. controls (-1.6,2.75) and (-1.4,2.75)
    .. (-1.5,2.25)
    .. controls (-1.6,1.75) and (-1.4,1.75)
    .. (-1.1,2.25);
  \node[left, codexblue] at (-2.2,2.25)
    {$V_p$ (AC)};

  % AC output
  \draw[thick, codexred] (6.5,4.5) -- (8.5,4.5);
  \draw[thick, codexred] (6.5,1.0) -- (8.5,1.0);
  % Load resistor
  \fill[gray!20] (8.2,1.5) rectangle (8.8,4.0);
  \draw[thick] (8.2,1.5) rectangle (8.8,4.0);
  \node[right, font=\small] at (8.9,2.75) {Load $R$};
  \draw[thick, codexred] (8.5,4.5) -- (8.5,4.0);
  \draw[thick, codexred] (8.5,1.0) -- (8.5,1.5);
  \node[right, codexred] at (9.1,4.0) {$V_s$};

  % Flux arrow
  \draw[->, very thick, codexgold]
    (0.25,0.75) -- (0.25,4.75);
  \draw[->, very thick, codexgold]
    (0.25,5.25) -- (6.75,5.25);
  \draw[->, very thick, codexgold]
    (6.75,4.75) -- (6.75,0.75);
  \draw[->, very thick, codexgold]
    (6.75,0.25) -- (0.25,0.25);
  \node[above, codexgold, font=\tiny]
    at (3.5,0.45) {Magnetic flux $\Phi$};

  % Turns ratio
  \node[align=center, font=\small]
    at (3.5,-1.0)
    {$\dfrac{V_s}{V_p} = \dfrac{N_s}{N_p}$
    \quad (turns ratio)};

\end{tikzpicture}
\end{center}
\end{diagbox}

\begin{lawbox}[Transformer Equations]
\textbf{Voltage ratio} (turns ratio):
\[
  \frac{V_s}{V_p} = \frac{N_s}{N_p}
\]

\textbf{Current ratio} (assuming 100\% efficiency):
\[
  \frac{I_s}{I_p} = \frac{N_p}{N_s}
\]
(Current ratio is the inverse of voltage ratio.)

\textbf{Power equality} (ideal transformer):
\[
  P_p = P_s \implies V_p I_p = V_s I_s
\]

\textbf{Step-up transformer:} $N_s > N_p$,
$V_s > V_p$, $I_s < I_p$.

\textbf{Step-down transformer:} $N_s < N_p$,
$V_s < V_p$, $I_s > I_p$.

\textbf{Efficiency:}
\[
  \eta = \frac{P_s}{P_p} \times 100\%
  = \frac{V_s I_s}{V_p I_p} \times 100\%
\]
\end{lawbox}

\begin{keybox}[Why Transformers Only Work with AC]
Transformers require a \textbf{changing} magnetic
flux to induce an EMF in the secondary coil. A DC
current through the primary creates a constant flux
--- no change, no induction, no output in the
secondary. Only AC (alternating current), which
reverses direction many times per second, creates
the continuously changing flux needed.

This is why AC was chosen over DC for power
distribution worldwide (following the famous
``War of Currents'' between Edison, who promoted
DC, and Tesla/Westinghouse, who promoted AC in
the 1880s--1890s). Tesla won.
\end{keybox}

\begin{examplebox}[24.5]
\stepcounter{workedexample}
A step-down transformer connected to a $240\ \text{V}$
AC supply has a turns ratio of $20:1$ and supplies
a resistive load of $6.0\ \Omega$.\\
(a) Find the secondary voltage.\\
(b) Find the secondary current.\\
(c) Find the primary current (assume 100\% efficiency).\\
(d) Find the power delivered to the load.\\
(e) The transformer efficiency is actually $95\%$.
Find the actual primary current.
\end{examplebox}

\begin{solutionbox}
\textbf{(a)} Secondary voltage:
\[
  V_s = V_p \times \frac{N_s}{N_p}
  = 240 \times \frac{1}{20} = 12\ \text{V}
\]

\textbf{(b)} Secondary current:
\[
  I_s = \frac{V_s}{R} = \frac{12}{6.0} = 2.0\ \text{A}
\]

\textbf{(c)} Primary current (ideal):
\[
  I_p = I_s \times \frac{N_s}{N_p}
  = 2.0 \times \frac{1}{20} = 0.10\ \text{A}
\]
Or: $I_p = P/(V_p) = 24/240 = 0.10\ \text{A}$ ✓

\textbf{(d)} Power to load:
\[
  P = I_s^2 R = (2.0)^2 \times 6.0 = 24\ \text{W}
\]
Or: $P = V_s I_s = 12 \times 2.0 = 24\ \text{W}$ ✓

\textbf{(e)} At $95\%$ efficiency:
\[
  \eta = \frac{P_s}{P_p}
  \implies P_p = \frac{P_s}{\eta}
  = \frac{24}{0.95} = 25.26\ \text{W}
\]
\[
  I_p = \frac{P_p}{V_p} = \frac{25.26}{240}
  = 0.105\ \text{A}
\]
The primary current is slightly higher than the
ideal case to account for the transformer's own
internal losses.
\end{solutionbox}

% ============================================================
\section{Alternating Current: Basic Concepts}
% ============================================================

\begin{definitionbox}[AC Quantities]
In \textbf{alternating current} (AC), the current
and voltage alternate sinusoidally:
\[
  V = V_0\sin(2\pi ft) = V_0\sin(\omega t)
\]
\[
  I = I_0\sin(\omega t)
\]

\textbf{Peak value:} $V_0$ or $I_0$ (maximum value).

\textbf{Root mean square (RMS) value:}
\[
  V_{rms} = \frac{V_0}{\sqrt{2}} \approx 0.707\,V_0
\]
\[
  I_{rms} = \frac{I_0}{\sqrt{2}} \approx 0.707\,I_0
\]

The RMS value of an AC current or voltage is the
equivalent DC value that produces the same power
dissipation.

\textbf{Power in AC circuit} (resistive load only):
\[
  P = V_{rms}I_{rms} = I_{rms}^2R
  = \frac{V_{rms}^2}{R}
\]
\end{definitionbox}

\begin{keybox}[Mains Voltage: Peak vs RMS]
The Nigerian (and UK) mains supply is quoted as
$240\ \text{V}$. This is the \textbf{RMS value}.
The actual peak voltage is:
\[
  V_0 = \sqrt{2} \times V_{rms}
  = \sqrt{2} \times 240 \approx 339\ \text{V}
\]
A voltmeter connected to the mains reads the RMS
value ($240\ \text{V}$). The peak value of $339\ \text{V}$
is high enough to be lethal --- which is why the
peak voltage matters for insulation design and the
rating of capacitors and other components.
\end{keybox}

\begin{examplebox}[24.6]
\stepcounter{workedexample}
The Nigerian mains supply is $240\ \text{V}$ RMS
at $50\ \text{Hz}$.\\
(a) Find the peak voltage.\\
(b) Find the peak current through a $60\ \Omega$
resistor connected to the mains.\\
(c) Find the average power dissipated.\\
(d) Write the equation for the instantaneous
voltage as a function of time.
\end{examplebox}

\begin{solutionbox}
\textbf{(a)} Peak voltage:
\[
  V_0 = \sqrt{2} \times 240 = 339\ \text{V}
\]

\textbf{(b)} Peak current:
\[
  I_0 = \frac{V_0}{R} = \frac{339}{60}
  = 5.65\ \text{A}
\]

\textbf{(c)} Average power (using RMS values):
\[
  P = \frac{V_{rms}^2}{R} = \frac{240^2}{60}
  = \frac{57\,600}{60} = 960\ \text{W}
\]

\textbf{(d)} Instantaneous voltage:
$\omega = 2\pi f = 2\pi \times 50 = 314\ \text{rad\,s}^{-1}$
\[
  V = 339\sin(314t)\ \text{V}
\]
\end{solutionbox}

% ============================================================
\section{Sources of Energy for Electricity Generation}
% ============================================================

\begin{table}[H]
\centering
\caption{Sources of Electrical Energy}
\label{tab:energy_sources}
\begin{tabular}{llll}
\toprule
\textbf{Source} & \textbf{Type} & \textbf{Nigerian
example} & \textbf{Efficiency}\\
\midrule
Coal/Gas & Fossil fuel & Egbin Gas Station & $35$--$45\%$\\
Hydroelectric & Renewable & Kainji, Shiroro & $85$--$95\%$\\
Nuclear & Nuclear fission & (not in Nigeria) & $33$--$37\%$\\
Solar PV & Renewable & Off-grid Nigeria & $15$--$22\%$\\
Wind & Renewable & (limited in Nigeria) & $35$--$45\%$\\
Geothermal & Renewable & (not applicable) & $10$--$23\%$\\
Biomass & Renewable & Sawmill waste & $20$--$30\%$\\
\bottomrule
\end{tabular}
\end{table}

\begin{realworldbox}[Nigeria's Electricity Challenges]
Nigeria has an installed electricity generation
capacity of approximately $12\,500\ \text{MW}$
(as of 2025), but typically generates only
$4\,000$--$5\,000\ \text{MW}$ against a demand
estimated at $30\,000\ \text{MW}$. The result is
the persistent power outages that have made private
generators and solar inverter systems a fact of
life for Nigerian homes and businesses.

Understanding electrical power, efficiency, and
energy costs is not abstract for Nigerians --- it
is a daily practical necessity. The physics in
this chapter directly governs the cost of running
a generator, the savings from switching to LED
bulbs, the economics of solar panels, and the
reason the national grid transmits at $330\ \text{kV}$
rather than $240\ \text{V}$.
\end{realworldbox}

% ============================================================
\section{Chapter Summary}
% ============================================================

\begin{summarybox}
\textbf{Electrical power:}
$P = IV = I^2R = V^2/R$ (watts, W)

\textbf{Electrical energy:}
$W = Pt$ (joules); $W = P(\text{kW}) \times t(\text{hours})$
(kWh); $1\ \text{kWh} = 3.6\ \text{MJ}$

\textbf{Joule heating:} $Q = I^2Rt$;
always generates heat regardless of current direction.

\textbf{Fuse:} melts when $I > $ rated value;
protects circuit. Choose rating just above $I_{operating}$.

\textbf{Power transmission:} at high voltage to
reduce $I$, reducing $P_{loss} = I^2R_{cable}$.
$P_{loss} \propto 1/V^2$ for same transmitted power.

\textbf{Transformer:}
$V_s/V_p = N_s/N_p$ (turns ratio);
$I_s/I_p = N_p/N_s$;
$P_p = P_s$ (ideal).
Requires AC. Step-up: $N_s > N_p$.
Step-down: $N_s < N_p$.

\textbf{AC:} $V_{rms} = V_0/\sqrt{2}$;
$I_{rms} = I_0/\sqrt{2}$.
Mains $240\ \text{V}$ is RMS; peak $= 339\ \text{V}$.
$P = V_{rms}I_{rms}$ (resistive load).
\end{summarybox}

% ============================================================
\newpage
\begin{exercisebox}[24]

\subsection*{Section A: Multiple Choice}

\textbf{A1.} A $100\ \Omega$ resistor carries a
current of $2.0\ \text{A}$. The power dissipated is:

\mcoption{A}{$50\ \text{W}$}
\mcoption{B}{$200\ \text{W}$}
\mcoption{C}{$400\ \text{W}$}
\mcoption{D}{$20\,000\ \text{W}$}
\ansref{Ch24-A1}

\vspace{0.4cm}

\textbf{A2.} $1\ \text{kWh}$ is equal to:

\mcoption{A}{$1000\ \text{J}$}
\mcoption{B}{$3600\ \text{J}$}
\mcoption{C}{$3.6\times10^6\ \text{J}$}
\mcoption{D}{$1\ \text{W}$}
\ansref{Ch24-A2}

\vspace{0.4cm}

\textbf{A3.} A transformer has a turns ratio
$N_p:N_s = 1:10$. It is a:

\mcoption{A}{Step-down transformer}
\mcoption{B}{Step-up transformer}
\mcoption{C}{Isolation transformer}
\mcoption{D}{Rectifier}
\ansref{Ch24-A3}

\vspace{0.4cm}

\textbf{A4.} Electrical power is transmitted at
high voltage in order to:

\mcoption{A}{Increase the speed of electricity}
\mcoption{B}{Reduce power losses in cables}
\mcoption{C}{Reduce the resistance of cables}
\mcoption{D}{Increase the efficiency of generators}
\ansref{Ch24-A4}

\vspace{0.4cm}

\textbf{A5.} The RMS value of an AC voltage with
peak value $100\ \text{V}$ is approximately:

\mcoption{A}{$50\ \text{V}$}
\mcoption{B}{$70.7\ \text{V}$}
\mcoption{C}{$100\ \text{V}$}
\mcoption{D}{$141.4\ \text{V}$}
\ansref{Ch24-A5}

\vspace{0.4cm}

\textbf{A6.} An ideal step-down transformer reduces
voltage from $240\ \text{V}$ to $12\ \text{V}$.
The ratio of primary current to secondary current is:

\mcoption{A}{$1:1$}
\mcoption{B}{$1:20$}
\mcoption{C}{$20:1$}
\mcoption{D}{$240:12$}
\ansref{Ch24-A6}

\vspace{0.4cm}

\textbf{A7.} A $2\ \text{kW}$ appliance is used for
$3\ \text{h}$ per day for 30 days. The energy consumed is:

\mcoption{A}{$6\ \text{kWh}$}
\mcoption{B}{$60\ \text{kWh}$}
\mcoption{C}{$180\ \text{kWh}$}
\mcoption{D}{$6000\ \text{kWh}$}
\ansref{Ch24-A7}

\vspace{0.4cm}

\textbf{A8.} A fuse blows when:

\mcoption{A}{The voltage across it exceeds the rating}
\mcoption{B}{The power input is too low}
\mcoption{C}{The current through it exceeds the
rated value and melts the wire}
\mcoption{D}{The resistance of the circuit is too high}
\ansref{Ch24-A8}

\vspace{0.4cm}

\textbf{A9.} If the transmission line voltage is
doubled (same power transmitted), the power
dissipated in the line:

\mcoption{A}{Doubles}
\mcoption{B}{Is halved}
\mcoption{C}{Is quartered}
\mcoption{D}{Stays the same}
\ansref{Ch24-A9}

\vspace{0.4cm}

\textbf{A10.} Transformers work only with AC because:

\mcoption{A}{AC has a higher voltage than DC}
\mcoption{B}{DC would destroy the iron core}
\mcoption{C}{Electromagnetic induction requires a
changing magnetic flux, which DC cannot provide}
\mcoption{D}{AC is safer than DC}
\ansref{Ch24-A10}

\vspace{0.6cm}

\subsection*{Section B: Theory and Calculations}

\textbf{B1.} A household in Abuja uses the
following electrical appliances:

\begin{center}
\begin{tabular}{lcc}
\toprule
Appliance & Power (W) & Hours per day\\
\midrule
Split air conditioner & 1800 & 8\\
Refrigerator & 150 & 24\\
LED bulbs (5 total) & 15 each & 6\\
Television & 120 & 5\\
Laptop charger & 65 & 4\\
Electric kettle & 1500 & 0.5\\
\bottomrule
\end{tabular}
\end{center}

The electricity tariff is ₦80 per kWh.\\[4pt]
(a) Calculate the daily energy consumption of
each appliance in kWh.\\
(b) Find the total daily energy consumption.\\
(c) Calculate the monthly electricity bill
(30-day month).\\
(d) The household switches from a $1800\ \text{W}$
split AC to a $1200\ \text{W}$ inverter AC with
the same cooling performance. Calculate the
monthly saving in Naira.\\
(e) Calculate the payback period (in months) if
the inverter AC costs ₦180,000 more than the
standard AC.
\hfill\ansref{Ch24-B1}

\vspace{0.4cm}

\textbf{B2.} A $220\ \text{kV}$ transmission line
carries power from the Kainji Hydroelectric Dam
to Lagos. The transmission line has a total
resistance (both conductors) of $60\ \Omega$.
The line carries $500\ \text{A}$.\\[4pt]
(a) Calculate the power delivered to the line.\\
(b) Calculate the power lost in the line as heat.\\
(c) Calculate the voltage drop along the line.\\
(d) What is the voltage at the Lagos end of the line?\\
(e) If the power delivered to Lagos consumers
is $P_L = (220\text{ kV} - \Delta V) \times 500\text{ A}$,
calculate $P_L$.\\
(f) Confirm the power balance:
$P_{input} = P_{line loss} + P_L$.
\hfill\ansref{Ch24-B2}

\begin{center}
\begin{tikzpicture}[scale=0.8,transform shape,font=\small,>=Stealth]
  \node[draw,minimum width=3.2cm,minimum height=1.0cm,align=center] (src) at (0,1.5) {Kainji station ($220\,\text{kV}$)};
  \node[draw,minimum width=2.5cm,minimum height=1.0cm,align=center] (load) at (8.5,1.5) {Lagos consumers};
  \draw[thick] (1.6,1.9) -- (7.25,1.9);
  \draw[thick] (7.25,1.1) -- (1.6,1.1);
  \draw[->,thick] (2.0,1.9) -- (3.0,1.9) node[midway,above] {$I=500\,\text{A}$};
  \draw[->,thick] (5.8,1.1) -- (4.8,1.1);
  \node[font=\footnotesize,fill=white] at (4.4,1.48) {$R_{\text{line}}=60\,\Omega$};
  \node[below,font=\footnotesize] at (4.4,0.65) {Total includes both conductors; find the voltage drop and power balance.};
\end{tikzpicture}
\end{center}

\vspace{0.4cm}

\textbf{B3.} A $240\ \text{V}$ to $12\ \text{V}$
step-down transformer supplies a car battery charger.
The secondary coil has $60$ turns. The charger
delivers $5.0\ \text{A}$ to the battery.\\[4pt]
(a) Find the number of primary turns.\\
(b) Find the primary current (assume 100\%
efficiency).\\
(c) Find the power input and output.\\
(d) The transformer has an efficiency of $92\%$.
Find the actual primary current.\\
(e) Calculate the heat generated inside the
transformer per hour.\\
(f) Explain two causes of energy loss in a
real transformer.
\hfill\ansref{Ch24-B3}

\vspace{0.4cm}

\textbf{B4.} An AC generator produces a sinusoidal
voltage with peak value $325\ \text{V}$ at
$50\ \text{Hz}$.\\[4pt]
(a) Write the expression for the instantaneous
voltage as a function of time.\\
(b) Calculate the RMS voltage.\\
(c) This generator is connected to a $40\ \Omega$
resistive load. Calculate the peak current, RMS
current, and average power.\\
(d) A voltmeter and an ammeter are connected to
measure this circuit. What do they read?\\
(e) If this generator is connected to the primary
of a transformer with $N_p = 200$ and $N_s = 50$,
find the secondary RMS voltage and the secondary
peak voltage.
\hfill\ansref{Ch24-B4}

\begin{center}
\begin{tikzpicture}
  \begin{axis}[
    width=10cm,height=4.2cm,
    xmin=0,xmax=0.04,ymin=-390,ymax=390,
    xtick={0,0.01,0.02,0.03,0.04},
    xticklabels={0,10,20,30,40},
    scaled x ticks=false,
    ytick={-325,0,325},
    xlabel={time (ms)},ylabel={instantaneous voltage (V)},
    axis lines=left,grid=major,
    minor tick num=1,
    tick label style={font=\footnotesize},
    label style={font=\small},
    clip=false
  ]
    \addplot[very thick,codexblue,domain=0:0.04,samples=160]
      {325*sin(deg(100*pi*x))};
  \end{axis}
\end{tikzpicture}
\par\footnotesize Reference waveform: peak $325\,\text{V}$, frequency $50\,\text{Hz}$; determine the requested RMS quantities yourself.
\end{center}

\vspace{0.4cm}

\textbf{B5.} A house is wired with cables rated
at a maximum current of $13\ \text{A}$ from a
$240\ \text{V}$ supply.\\[4pt]
(a) Calculate the maximum power that can be
drawn from each socket.\\
(b) An extension cord with four sockets is used.
All four sockets are used: $1800\ \text{W}$
air fryer, $900\ \text{W}$ microwave,
$1200\ \text{W}$ electric kettle, and a
$60\ \text{W}$ lamp. Calculate the total current.\\
(c) Will the $13\ \text{A}$ fuse blow? What is the
risk if the fuse is replaced with a $20\ \text{A}$
fuse?\\
(d) The air fryer, microwave and kettle each have
internal resistance of approximately
$32\ \Omega$, $64\ \Omega$ and $48\ \Omega$
respectively. Calculate the power dissipated in the
extension cord itself if the cord has total
resistance $0.5\ \Omega$.\\
(e) A student argues that if the voltage is $240\ \text{V}$
and the fuse is rated $13\ \text{A}$, you can safely
use up to $240 \times 13 = 3120\ \text{W}$. Explain
whether this is correct, and identify any assumptions.
\hfill\ansref{Ch24-B5}

\begin{center}
\begin{tikzpicture}[scale=0.72,transform shape,font=\small,>=Stealth]
  % The fuse and cord resistance are in series with the shared live feed.
  \draw[thick] (0.5,0.45) -- (0.5,2.3);
  \draw[thick] (0.5,3.25) -- (0.5,3.23);
  \node[draw,minimum width=7mm,minimum height=3mm,inner sep=1pt] at (0.5,3.08) {$F$};
  \draw[thick] (0.5,2.93) -- (0.5,2.79);
  \node[draw,minimum width=7mm,minimum height=3mm,inner sep=1pt] at (0.5,2.64) {$R_c$};
  \draw[thick] (0.5,2.49) -- (0.5,2.3);
  \draw[thick] (7.0,0.45) -- (7.0,3.25);
  \node[left] at (0.5,3.45) {$240\,\text{V}$};
  \node[anchor=west,font=\scriptsize] at (0.95,3.08) {$13\,\text{A}$ fuse};
  \node[anchor=west,font=\scriptsize] at (0.95,2.64) {cord $R_c=0.5\,\Omega$};
  \node[right] at (7.0,3.0) {return};
  \draw[->,thick] (1.4,3.45) -- (2.3,3.45) node[midway,above] {$I_{total}$};
  \foreach \y/\name/\power in {2.3/Air fryer/1800 W,1.75/Microwave/900 W,1.2/Kettle/1200 W,0.65/Lamp/60 W} {
    \draw[thick] (0.5,\y) -- (2.05,\y);
    \draw[thick] (2.05,\y-0.22) rectangle (4.75,\y+0.22);
    \node at (3.4,\y) {\name\space(\power)};
    \draw[thick] (4.75,\y) -- (7.0,\y);
  }
  \node[below] at (3.75,0.12) {Four socket branches are in parallel and share the fused cord.};
\end{tikzpicture}
\end{center}

\end{exercisebox}
